%!TEX root = TQFTbook.tex
%%%%%%%%%%%%%%%%%%%%%%%%%%%%%%
%\label{c:p1-overview}
In this part, we introduce the basic definitions
of TQFT.

In \chref{c:physical-motivation}, we briefly review the main ideas of quantum
field theory in physics, most importantly path integral formulation of QFT due
to Feynman. This section is extremely brief and is used as motivation only; no
precise definitions or statements are given there. However, we hope that it
helps the reader understand the physical roots of TQFT.

In \chref{c:categories}, we introduce the key mathematical tool which will be
extensively used throughout this book: the notions of a monoidal category,
monoidal functor, and related topics.

Finally, in \chref{c:cobordisms} we define the category of cobordisms and use
it to give a definition of TQFT as a functor $\nCob\to \vect$. This definition
is essentially due to Atiyah. We also study some basic properties of TQFTs, in
particular showing that a TQFT $Z$ defines, for any closed $(n-1)$-manifold
$N$, an action of the mapping class group $\Ga(N)=\pi_0(\Diff_+(N))$ on the
vector space $Z(N)$. As a baby example, we give a classification of
1-dimensional TQFTs. This is easy, but serves as a foundation for all
classification results which appear later in the book, such as the cobordism
hypothesis.

In \chref{c:2d-classification}, we discuss classification of 2-dimensional TQFT,
showing that such theories are classified by commutative Frobenius algebras. Our
approach is based on the use of Morse functions and families of Morse functions
(Cerf theory).

In \chref{c:frobenius}, we study Frobenius algebras, in particular classifying
all semisimple commutative Frobenius algebras and describing the corresponding
2d TQFTs.

In \chref{c:DW}, we describe the fundamental example of a TQFT: gauge theory
with finite gauge group, also known as Dijkgraaf--Witten theory. This theory can
be described in terms of $G$-covers of surfaces; we can also relate it to the
path integral description of TQFT.
